% ===========================================================================
% Poste 9 — Crédit d'impôt pour la solidarité
% ===========================================================================
\poste{9}{Crédit d'impôt pour la solidarité}{QC\_sol}

\pdesc Crédit remboursable versé mensuellement par Revenu Québec, qui
regroupe depuis 2011 trois anciennes mesures : la compensation de la TVQ, le
remboursement d'impôts fonciers (volet logement) et le soutien aux
habitants des villages nordiques.

\pobj Compenser la taxe de vente du Québec et une partie des coûts de
logement des ménages à revenu faible ou moyen.

\pregle Le crédit additionne deux volets, puis applique une réduction unique
de \pc{6} sur le revenu familial net excédant le seuil $s$ :
\begin{align*}
\text{volet TVQ} &= B + \begin{cases}
C & \text{si couple (montant pour conjoint)} \\
A & \text{si un seul adulte (montant \og vivant seule \fg{})}
\end{cases} \\
\text{volet logement} &= \begin{cases}
L_{couple} & \text{si couple} \\
L_{seul} & \text{sinon}
\end{cases} \;+\; n \times L_{enfant} \\
\text{crédit} &= \pos{\text{volet TVQ} + \text{volet logement}
 - 0{,}06 \times \pos{\RFN - s}}
\end{align*}
Le modèle suppose le ménage locataire ou propriétaire d'un logement
admissible (droit aux deux volets) ; le taux de réduction est donc toujours
\pc{6} (il serait de \pc{3} à une seule composante). Le volet
\og village nordique \fg{} n'est pas modélisé. La famille monoparentale
reçoit le montant additionnel \og vivant seule \fg{}, étant réputée vivre
seule avec ses enfants à charge. Les périodes de versement vont de juillet à
juin ; l'année affichée correspond à la période débutant en juillet.

\pparams

\begin{center}
\small
\begin{tabular}{l r r}
\toprule
Paramètre & 2025 & 2026 \\
\midrule
Volet TVQ --- montant de base ($B$)            & \D{356} & \D{363} \\
Volet TVQ --- montant pour conjoint ($C$)      & \D{356} & \D{363} \\
Volet TVQ --- additionnel vivant seule ($A$)   & \D{169} & \D{172} \\
Volet logement --- couple ($L_{couple}$)       & \D{888} & \D{906} \\
Volet logement --- seule ou monoparentale ($L_{seul}$) & \D{731} & \D{746} \\
Volet logement --- par enfant ($L_{enfant}$)   & \D{155} & \D{158} \\
Seuil de réduction ($s$)                       & \D{42325} & \D{43195} \\
Taux de réduction                              & \pc{6}  & \pc{6} \\
\bottomrule
\end{tabular}
\end{center}

\pnotes Le volet logement exige en réalité un logement \emph{admissible}
--- locataire d'un logement visé par un relevé~31 ou propriétaire ; sont
notamment exclus les habitations à loyer modique et les logements dont le
loyer est subventionné --- condition que le modèle tient pour acquise. La
\emph{garde partagée} (demi-montant par enfant au volet logement) n'est pas
couverte par les entrées. Enfin, les prestataires d'une aide de dernier
recours reçoivent le crédit automatiquement, sans demande.

% ============================================================================
% GÉNÉRÉ par scripts/exemples-guide.ts — NE PAS ÉDITER À LA MAIN.
% Régénérer : npm run exemples (chaque chiffre est vérifié contre le moteur).
% ============================================================================
\pexemples Deux volets (TVQ et logement) puis une réduction de \pc{6}
du $\RFN$ au-delà de \D{42325}.

\medskip\noindent\textbf{M3 --- personne vivant seule, 30~ans, revenu de travail de \D{15000}.}
\begin{align*}
\text{volet TVQ} &= \num{356} + \num{169} = \num{525} \quad (\text{base} + \text{vivant seul}) \\
\text{volet logement} &= \num{731} \\
\text{réduction} &= 0 \quad (\RFN = \num{13985} < \num{42325})
\end{align*}
Crédit \emph{plein} : \D{1256}.

\medskip\noindent\textbf{M4 --- personne vivant seule, 40~ans, revenu de travail de \D{50000}.}
Mêmes volets (\num{525} et \num{731}), mais le $\RFN$
excède le seuil :
\begin{align*}
\text{réduction} &= \pos{\num{48115} - \num{42325}} \times \pc{6} = \num{347.40} \\
\text{crédit} &= \pos{\num{1256} - \num{347.40}} = \num{908.60}
\end{align*}
Crédit \emph{partiel} : \D{908.60}.

\medskip\noindent\textbf{M8 --- couple, 38 et 36~ans, revenus de travail de \D{60000} et \D{40000}, deux enfants : 4~ans (garde non subventionnée, \D{13000}) et 8~ans (garde non subventionnée, \D{3000}).}
Couple avec deux enfants :
\begin{align*}
\text{volet TVQ} &= \num{356} + \num{356} = \num{712} \\
\text{volet logement} &= \num{888} + 2 \times \num{155} = \num{1198} \\
\text{réduction} &= \pos{\num{96230} - \num{42325}} \times \pc{6} = \num{3234.30} \\
\text{crédit} &= \pos{\num{1910} - \num{3234.30}} = \num{0}
\end{align*}
Crédit : \D{0}.

\medskip\noindent\textbf{M5 --- personne vivant seule, 45~ans, revenu de travail de \D{100000}.}
La réduction (\pc{6} de \num{97506} $-$
\num{42325} $=$ \num{3310.86}) dépasse la somme des volets
(\num{1256}) : crédit \emph{éteint} (\D{0}).


\prefs Loi sur les impôts (RLRQ, c.~I-3), art.~1029.8.116.12 à
1029.8.116.35 ; Revenu Québec ; ministère des Finances du Québec,
\emph{Paramètres du régime d'imposition} (tableau du crédit pour la
solidarité) ; CFFP, \emph{Guide des mesures fiscales}, fiche \og Crédit
d'impôt pour la solidarité \fg{}\footnote{\url{https://cffp.recherche.usherbrooke.ca/outils-ressources/guide-mesures-fiscales/credit-impot-solidarite/}}.
